\documentclass{beamer}
\setbeamertemplate{navigation symbols}{}
\setbeamertemplate{footline}[frame number]
% \usepackage{mathpazo}

\usepackage{graphicx}
\usepackage{amsmath}
\graphicspath{ {../pictures/} } 

\title{3. Variations of Brownian Motion}
\date{\today}

\begin{document}
	
	\begin{frame}
		\titlepage
		\centering All proofs from Stochastic Processes by Sheldon Ross
	\end{frame}
	
	% ---------------------------------------------------
	% 3.1 Absorbed Brownian Motion
	% ---------------------------------------------------
	\begin{frame}{3.1: Absorbed Brownian Motion}
		An absorbed Brownian motion behaves exactly like a standard Brownian motion until it hits a specific barrier $a$. Once it reaches $a$, the process is absorbed and remains fixed at that value forever.
		
		\vspace{0.5cm}
		\textbf{Definition:} Let $\tau_a = \inf\{t \ge 0 : W(t) = a\}$. The absorbed process $X(t)$ is:
		$$X(t) = \begin{cases}
			W(t) & \text{if } t < \tau_a \\
			a & \text{if } t \ge \tau_a,
		\end{cases}$$
		
	\end{frame}
	
	\begin{frame}{3.1: Absorbed Brownian Motion (Simulation)}
		\begin{figure}
			\centering
			\includegraphics[width=\textwidth]{"2_absorbed_brownian.png"}
			\caption{Paths and distribution of Absorbed Brownian Motion at barrier $a=0.5$}
		\end{figure}
	\end{frame}
	
	\begin{frame}{3.1: Derivation of Absorbed BM PDF}
		Let $W(t)$ be a standard Brownian motion and $a > 0$ be an absorbing barrier.
		The stopping time is
		$$\tau_a = \inf\{t > 0 : W(t) = a\}$$.
		
		For the absorbed process $X(t)$, we want to find the CDF for $x < a$:
		$$P(X(t) \le x) = P(W(t) \le x, \tau_a > t)$$
		
		Using the Law of Total Probability:
		$$P(W(t) \le x) = P(W(t) \le x, \tau_a > t) + P(W(t) \le x, \tau_a \le t)$$
		
		By symmetry, in standard brownian motion, paths can hit $x<a$ or $a+(a-x) = 2a-x$ at time $t \ge \tau_a$ with equal probability (going up or down the same amount has equal probability, this amount is $a-x$).
		$$P(W(t) \le x, \tau_a \le t) = P(W(t) \ge 2a - x, \tau_a \le t)$$
		
	\end{frame}
	
	\begin{frame}{3.1: Derivation of Absorbed BM PDF}
		Now observe that $W(t) > 2a-x$ implies $W(t) > a$ (since $x<a$). Thus, this means that $t \ge \tau_a$ and so the probability can be written as:
		$$P(W(t) \ge 2a - x, \tau_a \le t) = P(W(t) \ge 2a - x)$$
		
		Substituting the probability back into the CDF equation:
		$$P(X(t) \le x) = P(W(t) \le x) - P(W(t) \ge 2a - x)$$
		
		Using the standard normal CDF $\Phi$:
		$$F_{X(t)}(x) = \Phi\left(\frac{x}{\sqrt{t}}\right) - \left(1 - \Phi\left(\frac{2a - x}{\sqrt{t}}\right)\right)$$
		
		Since the normal distribution is symmetric, $1 - \Phi(z) = \Phi(-z)$:
		$$F_{X(t)}(x) = \Phi\left(\frac{x}{\sqrt{t}}\right) - \Phi\left(\frac{x - 2a}{\sqrt{t}}\right)$$
		
		Differentiating with respect to $x$ gives the PDF for $x < a$:
		$$f_{X(t)}(x) = \frac{1}{\sqrt{2\pi t}} \left( e^{-\frac{x^2}{2t}} - e^{-\frac{(x - 2a)^2}{2t}} \right)$$
	\end{frame}
	
	\begin{frame}{3.1: Hitting Time Probability in Standard BM}
		Now consider standard Brownian motion $W(t)$. The CDF of the hitting time $\tau_a$ is $P(\tau_a \le t)$.
		$$P(\tau_a \le t) = P(W(t) \ge a, \tau_a \le t) + P(W(t) < a, \tau_a \le t)$$
		
		By symmetry, $P(W(t) < a, \tau_a \le t) = P(W(t) \ge a, \tau_a \le t) = P(W(t) \ge a)$. Thus:
		$$P(\tau_a \le t) = 2P(W(t) \ge a) = 2 \left( 1 - \Phi\left(\frac{a}{\sqrt{t}}\right) \right)$$
		
		To find the probability of hitting $a$ in finite time, we take the limit as $t \to \infty$:
		$$\lim_{t \to \infty} P(\tau_a \le t) = 2 \left( 1 - \Phi(0) \right) = 2(1 - 0.5) = 1$$
		
		Therefore, standard Brownian motion hits any value $a$ in finite time with probability $1$.
	\end{frame}
	
	\begin{frame}{3.1: Expected Hitting Time is Infinite}
		To find the expected hitting time $\mathbb{E}[\tau_a]$, we first find the PDF of $\tau_a$ by differentiating its CDF with respect to $t$:
		$$f_{\tau_a}(t) = \frac{d}{dt} \left[ 2 \int_{a/\sqrt{t}}^{\infty} \frac{1}{\sqrt{2\pi}} e^{-\frac{z^2}{2}} dz \right] = \frac{a}{\sqrt{2\pi t^3}} e^{-\frac{a^2}{2t}}$$
		
		As $t \to \infty$, the integrand behaves like $t^{-1/2}$, which diverges (proof in next slide). Hence:
		$$\mathbb{E}[\tau_a] = \infty$$
		
		Because $P(\tau_a < \infty) = 1$ and $\mathbb{E}[\tau_a] = \infty$ for any arbitrary state $a$, every state in a standard Brownian motion is \textbf{null recurrent}.
	\end{frame}
	
	
	\begin{frame}{3.1: Expected Hitting Time is Infinite}
		The expected value is:
		\begin{align*}
			\mathbb{E}[\tau_a] &= \int_0^\infty P(\tau_a > t) \, dt \\
			&= \int_0^\infty \left( 1 - \frac{2}{\sqrt{2\pi}} \int_{a/\sqrt{t}}^\infty e^{-y^2/2} \, dy \right) dt \\
			&= \frac{2}{\sqrt{2\pi}} \int_0^\infty \int_0^{a/\sqrt{t}} e^{-y^2/2} \, dy \, dt \\
			&= \frac{2}{\sqrt{2\pi}} \int_0^\infty \int_0^{a^2/y^2} dt \, e^{-y^2/2} \, dy \\
			&= \frac{2a^2}{\sqrt{2\pi}} \int_0^\infty \frac{1}{y^2} e^{-y^2/2} \, dy \\
			&\ge \frac{2a^2 e^{-1/2}}{\sqrt{2\pi}} \int_0^1 \frac{1}{y^2} \, dy \\
			&= \infty.
		\end{align*}
	\end{frame}
	
	% ---------------------------------------------------
	% 3.2 Reflected Brownian Motion
	% ---------------------------------------------------
	\begin{frame}{3.2: Reflected Brownian Motion at Origin}
		Reflected Brownian motion behaves like a standard Brownian motion but is constrained to be non-negative. Whenever the process hits $0$, it bounces back into the positive domain.
		
		\vspace{0.5cm}
		\textbf{Definition:} The standard reflected Brownian motion is the absolute value of a standard Brownian motion:
		$$Z(t) = |W(t)|$$
	\end{frame}
	
	\begin{frame}{3.2: Reflected Brownian Motion (Simulation)}
		\begin{figure}
			\centering
			\includegraphics[width=\textwidth]{"3_reflected_brownian.png"}
			\caption{Paths and distribution of Reflected Brownian Motion}
		\end{figure}
	\end{frame}
	
	\begin{frame}{3.2: Distribution of Reflected Brownian Motion}		
		The distribution of $Z(t)$ is obtained as:
		
		For $y > 0$,
		\begin{align*}
			P(Z(t) \le y) &= P(-y \le W(t) \le y) \\
			&= P(W(t) \le y) - P(W(t) \le -y) \\
			&= 2P(W(t) \le y) - 1 \\
			&= \frac{2}{\sqrt{2\pi t}} \int_{-\infty}^{y} e^{-x^2 / 2t} \, dx - 1,
		\end{align*}
		($P(W(t) \le -y) = 1-P(W(t) \le y)$, due to W(t) having gaussian distribution)
		
		\begin{itemize}
			\item Expected value: $ \mathbb{E}[Z(t)] = \sqrt{\frac{2t}{\pi}} $
			\item Variance: $ \text{Var}(Z(t)) = \left( 1 - \frac{2}{\pi} \right) t $
		\end{itemize}
		
	\end{frame}
	
	% ---------------------------------------------------
	% 3.3 Geometric Brownian Motion
	% ---------------------------------------------------
	\begin{frame}{3.3: Geometric Brownian Motion}
		If $\{W(t), t \ge 0\}$ is a standard Brownian motion, then the process $\{Y(t), t \ge 0\}$, defined by
		$$Y(t) = e^{W(t)}, \quad t \ge 0,$$
		is called \textbf{geometric Brownian motion}.
	\end{frame}
	
	\begin{frame}{3.3: Geometric Brownian Motion (Simulation)}
		\begin{figure}
			\centering
			\includegraphics[width=\textwidth]{"4_geometric_brownian.png"}
			\caption{Paths and Log-Normal distribution of GBM}
		\end{figure}
	\end{frame}
	
	\begin{frame}{3.3: Geometric Brownian Motion: Definition and Moments}
		Since $W(t)$ is normal with mean $0$ and variance $t$, its MGF is given by:
		$$\mathbb{E}[e^{sW(t)}] = e^{ts^2/2}$$
		
		Using this, we can compute the mean and variance of Y(t):
		\begin{align*}
			\mathbb{E}[Y(t)] &= \mathbb{E}[e^{W(t)}] = e^{t/2} \\
			\text{Var}(Y(t)) &= \mathbb{E}[Y^2(t)] - (\mathbb{E}[Y(t)])^2 \\
			&= \mathbb{E}[e^{2W(t)}] - e^t \\
			&= e^{2t} - e^t.
		\end{align*}
	\end{frame}
		
	\begin{frame}{3.3: GBM Example: Modelling stock prices}
		Using discrete real-world observations and assumptions (like stock prices):
		\begin{itemize}
			\item Let $Y(n)$ be a stock price at day $n$.
			\item Assume the daily \textbf{percentage returns}, $X_i = \frac{Y(i)}{Y(i-1)}$, are i.i.d.
		\end{itemize}
		
		\vspace{0.3cm}
		The price at time $n$ is the product of all previous returns:
		$$Y(n) = Y(0) X_1 X_2 \cdots X_n$$
		
		Taking the logarithm turns the product into a sum of i.i.d. variables:
		$$\log Y(n) = \log Y(0) + \sum_{i=1}^n \log X_i$$
		
		\vspace{0.3cm}
		By the \textbf{Central Limit Theorem}, this sum of i.i.d. variables converges to a Brownian motion as the time steps become continuous. Thus, $Y(t)$ naturally becomes a Geometric Brownian Motion.
	\end{frame}
	
	\begin{frame}{3.3: Discrete Prices and Multiplicative Returns}
		Before jumping to continuous calculus, we observe real-world markets in discrete time steps (e.g., daily closing prices).
		
		\begin{itemize}
			\item Let $Y(n)$ be the price of an asset at day $n$.
			\item Let $X_i = \frac{Y(i)}{Y(i-1)}$ be the gross return on day $i$. 
			\item We assume these daily returns $X_i$ are \textbf{independent and identically distributed (i.i.d.)}.
		\end{itemize}
		
		The price at day $n$ is simply the initial price multiplied by all subsequent daily returns:
		$$Y(n) = Y(0) X_1 X_2 \cdots X_n$$
	\end{frame}
	
	\begin{frame}{3.3: Discrete Prices and Multiplicative Returns}
		Taking the natural logarithm transforms the product into a sum:
		$$\log Y(n) = \log Y(0) + \sum_{i=1}^n \log X_i$$
		
		\begin{itemize}
			\item Let $Z_i = \log X_i$. Since the $X_i$ are i.i.d., the $Z_i$ terms are also i.i.d.
			\item Let the mean be $\mathbb{E}[Z_i] = \mu$ and variance be $\text{Var}(Z_i) = \sigma^2$.
			\item We define our sum: $S_n = \sum_{i=1}^n Z_i$.
		\end{itemize}
		
		The Central Limit Theorem (CLT) applies to our sum $S_n$, by appropriate scaling.
		
		Subtract its mean ($n\mu$) and divide by its standard deviation ($\sigma\sqrt{n}$):
		
		\begin{equation*}
			\frac{S_n - n\mu}{\sigma \sqrt{n}} \xrightarrow{d} \mathcal{N}(0, 1) \quad \text{as } n \to \infty
		\end{equation*}
		
		Thus $S_n$ becomes a gaussian random variable with mean $n\mu$ and variance $n\sigma^2$.
		$$S_n \sim \mathcal{N}(n\mu, n\sigma^2)$$
	\end{frame}

	% ---------------------------------------------------
	% 3.4 Integrated Brownian Motion
	% ---------------------------------------------------
	\begin{frame}{3.4: Integrated Brownian Motion}
		If $\{W(t), t \ge 0\}$ is a standard Brownian motion, the \textbf{integrated Brownian motion} $\{Y(t), t \ge 0\}$ is defined as the area under the Brownian path:
		$$Y(t) = \int_0^t W(s) \, ds$$
		
		\textbf{Physical Intuition:}
		\begin{itemize}
			\item If $W(t)$ represents the random \textit{velocity} of a particle subject to continuous random shocks...
			\item ...then $Y(t)$ represents the \textit{position} of that particle.
		\end{itemize}
		
		\textbf{Smoothness:} While $W(t)$ is jagged and nowhere differentiable, integration is a smoothing operation. $Y(t)$ is continuously differentiable, and its derivative is exactly $\frac{d}{dt}Y(t) = W(t)$. The path for IBM is smooth.
	\end{frame}
	
	\begin{frame}{3.4: Integrated Brownian Motion (Simulation)}
		\begin{figure}
			\centering
			\includegraphics[width=\textwidth]{"5_integrated_brownian.png"}
			\caption{Paths and distribution of Integrated Brownian Motion}
		\end{figure}
	\end{frame}
	
	\begin{frame}{3.4: Y(t) is Gaussian}
		The variables of a Brownian motion at any finite set of times are \textbf{jointly Gaussian}.
		
		\vspace{0.3cm}
		Consider any sequence of times $0 < t_1 < t_2 < \dots < t_n$. By the definition of Brownian motion, the \textit{increments} are independent and normally distributed:
		\begin{align*}
			Z_1 &= W(t_1) - W(0) \sim \mathcal{N}(0, t_1) \\
			Z_2 &= W(t_2) - W(t_1) \sim \mathcal{N}(0, t_2 - t_1) \\
			&\vdots \\
			Z_n &= W(t_n) - W(t_{n-1}) \sim \mathcal{N}(0, t_n - t_{n-1})
		\end{align*}
		
		Because the increments $Z_i$ are independent Gaussians, the vector $\mathbf{Z} = [Z_1, Z_2, \dots, Z_n]^T$ is trivially a \textbf{jointly Gaussian vector}.
	\end{frame}
	
	\begin{frame}{3.4: Y(t) is Gaussian}
		We can express the actual Brownian motion values $W(t_k)$ as a cumulative sum of these independent Gaussian increments:
		\begin{align*}
			W(t_1) &= Z_1 \\
			W(t_2) &= Z_1 + Z_2 \\
			&\vdots \\
			W(t_n) &= Z_1 + Z_2 + \dots + Z_n
		\end{align*}
		
		\vspace{0.3cm}
		\textbf{Property:} Any linear transformation of a jointly Gaussian vector results in another jointly Gaussian vector. 
		
		\vspace{0.3cm}
		Since $\mathbf{W} = [W(t_1), \dots, W(t_n)]^T$ is just a linear combination of $\mathbf{Z}$, we conclude that \textbf{$W(t)$ at any set of time instants is jointly Gaussian}.
	\end{frame}
	
	\begin{frame}{3.4: Y(t) is Gaussian}		
		Using the definition of Riemann integral as a limit of sums over small time steps $\Delta s$:
		$$Y(t) \approx \sum_{i=1}^n W(s_i) \Delta s$$
		
		\begin{itemize}
			\item From our previous proof, the variables $(W(s_1), W(s_2), \dots, W(s_n))$ are jointly Gaussian.
			\item A finite linear combination of jointly Gaussian variables is strictly Gaussian.
			\item The limit of a sequence of Gaussian random variables (as $\Delta s \to 0$) is also Gaussian.
		\end{itemize}
		
		Therefore, $Y(t)$ is a normally distributed random variable. We just need its parameters.
	\end{frame}
	
	\begin{frame}{3.4: Parameters of IBM: Expected Value}
		Since integration is a linear operator, we can pass the expectation directly inside the integral, indirectly applying linearity of expectation to variables in the Riemann integral
		
		The mean of the integrated Brownian motion is:
		\begin{align*}
			\mathbb{E}[Y(t)] &= \mathbb{E}\left[ \int_0^t W(s) \, ds \right] \\
			&= \int_0^t \mathbb{E}[W(s)] \, ds
		\end{align*}
		
		Because standard Brownian motion has a mean of zero at all times ($\mathbb{E}[W(s)] = 0$):
		$$ \mathbb{E}[Y(t)] = \int_0^t 0 \, ds = 0 $$
	\end{frame}
	
	\begin{frame}{3.4: Parameters of IBM: Variance}
		To find the variance, we calculate $\mathbb{E}[Y(t)^2]$ since the mean is zero. We square the integral by writing it as a double integral with dummy variables $u$ and $v$:
		
		\begin{align*}
			\text{Var}(Y(t)) &= \mathbb{E}\left[ \left( \int_0^t W(s) \, ds \right)^2 \right] \\
			&= \mathbb{E}\left[ \int_0^t W(u) \, du \int_0^t W(v) \, dv \right] \\
			&= \mathbb{E}\left[ \int_0^t \int_0^t W(u)W(v) \, du \, dv \right]
		\end{align*}
		
		\vspace{0.3cm}
		Passing the expectation inside the double integral:
		$$ \text{Var}(Y(t)) = \int_0^t \int_0^t \mathbb{E}[W(u)W(v)] \, du \, dv $$
	\end{frame}
	
	\begin{frame}{3.4: Covariance of Standard Brownian Motion}
		To evaluate the double integral for the variance, we need the covariance of standard Brownian motion: $\mathbb{E}[W(u)W(v)]$. Let's derive it.
		
		\vspace{0.3cm}
		Assume without loss of generality that $u \le v$. We can rewrite the value of the Brownian motion at the later time $v$ as its value at the earlier time $u$ plus the increment that occurred between $u$ and $v$:
		$$ W(v) = W(u) + (W(v) - W(u)) $$
		
		\vspace{0.3cm}
		Substitute this expanded form of $W(v)$ into the expectation:
		\begin{align*}
			\mathbb{E}[W(u)W(v)] &= \mathbb{E}\big[W(u) \cdot \big(W(u) + W(v) - W(u)\big)\big] \\
			&= \mathbb{E}\big[W(u)^2 + W(u)(W(v) - W(u))\big]
		\end{align*}
	\end{frame}
	
	\begin{frame}{3.4: Covariance of Standard Brownian Motion}
		Using the linearity of expectation, we split the expression into two terms, and using W(0) = 0:
		$$ \mathbb{E}[W(u)W(v)] = \mathbb{E}[W(u)^2] + \mathbb{E}[(W(u) - W(0))(W(v) - W(u))] $$
		
		\vspace{0.3cm}
		Now, we analyze the two terms:
		\begin{enumerate}
			\item \textbf{First Term:} Since $W(u) \sim \mathcal{N}(0, u)$, its second moment (which is its variance) is exactly $\mathbb{E}[W(u)^2] = u$.
			\item \textbf{Second Term:} Because Brownian motion has \textit{independent increments}, the increment $W(u) - W(0)$ is independent of the non overlapping future increment $W(v) - W(u)$. Thus:
			$$ \mathbb{E}[W(u) - W(0)] \cdot \mathbb{E}[W(v) - W(u)] = 0 \cdot 0 = 0 $$
		\end{enumerate}
		
		\vspace{0.3cm}
		Adding them yields $\mathbb{E}[W(u)W(v)] = u + 0 = u$. Similarly, if $v \le u$ instead, the result would be $v$. Therefore, the covariance is always $\min(u, v)$.
	\end{frame}
	
	\begin{frame}{3.4: Parameters of IBM: Covariance}
		Since $\{Y(t), t \ge 0\}$ is a Gaussian process, its distribution is fully characterized by its mean value and covariance function. 
		
		\vspace{0.3cm}
		For $s \le t$, we compute the covariance:
		\begin{align*}
			\text{Cov}[Y(s), Y(t)] &= \mathbb{E}[Y(s)Y(t)] \\
			&= \mathbb{E}\left[ \int_0^s W(y) \, dy \int_0^t W(u) \, du \right] \\
			&= \int_0^s \int_0^t \mathbb{E}[W(y)W(u)] \, du \, dy \\
			&= \int_0^s \int_0^t \min(y, u) \, du \, dy
		\end{align*}
	\end{frame}
	
	\begin{frame}{3.4: Parameters of IBM: Covariance}
		To evaluate $\int_0^s \int_0^t \min(y, u) \, du \, dy$ for $s \le t$, we split the inner integral at $u=y$:
		
		\begin{align*}
			\text{Cov}[Y(s), Y(t)] &= \int_0^s \left( \int_0^y u \, du + \int_y^t y \, du \right) dy \\
			&= \int_0^s \left( \frac{y^2}{2} + y(t - y) \right) dy \\
			&= \int_0^s \left( yt - \frac{y^2}{2} \right) dy \\
			&= t\frac{s^2}{2} - \frac{s^3}{6} \\
			&= s^2 \left( \frac{t}{2} - \frac{s}{6} \right)
		\end{align*}
		
		\vspace{0.2cm}
		\textit{Note: Setting $s=t$ immediately recovers the variance $t^3/3$.}
		
		\vspace{0.2cm}
		\textbf{Conclusion:} For any time $t$, $Y(t) \sim \mathcal{N}\left(0, \frac{t^3}{3}\right)$.
	\end{frame}
	
	\begin{frame}{3.4: Markov Process (Y(t), W(t))}
		While $Y(t)$ alone is not a Markov process, the vector process $\{(Y(t), W(t)), t \ge 0\}$ is Markov. Because they are jointly normal, their joint distribution is defined by their cross-covariance.
		
		\vspace{0.2cm}
		We analyze the difference over a small step $h$:
		\begin{align*}
			\text{Cov}(Y(t), Y(t) - Y(t-h)) &= \text{Cov}(Y(t), Y(t)) - \text{Cov}(Y(t), Y(t-h)) \\
			&= \frac{t^3}{3} - (t-h)^2 \left[ \frac{t}{2} - \frac{t-h}{6} \right] \\
			&= \frac{t^2h}{2} + o(h)
		\end{align*}
		
		\vspace{0.2cm}
		Alternatively, expressing the difference as an integral:
		\begin{align*}
			\text{Cov}(Y(t), Y(t) - Y(t-h)) &= \text{Cov}\left(Y(t), \int_{t-h}^t W(s) \, ds\right) \\
			&= \text{Cov}(Y(t), hW(t) + o(h)) \\
			&= h\text{Cov}(Y(t), W(t)) + o(h)
		\end{align*}
	\end{frame}
	
	\begin{frame}{3.4: The Bivariate Normal Distribution}
		Equating the two expressions for the covariance over step $h$:
		$$h\text{Cov}(Y(t), W(t)) \approx \frac{t^2h}{2}$$
		
		\vspace{0.3cm}
		Dividing by $h$ and taking the limit as $h \to 0$ gives the cross-covariance:
		$$\text{Cov}(Y(t), W(t)) = \frac{t^2}{2}$$
		
		\vspace{0.3cm}
		Hence, the joint state of the integrated process and the Brownian motion, $(Y(t), W(t))$, has a bivariate normal distribution with:
		\begin{align*}
			\mathbb{E}[Y(t)] &= \mathbb{E}[W(t)] = 0 \\
			\text{Var}(Y(t)) &= \frac{t^3}{3}, \quad \text{Var}(W(t)) = t \\
			\text{Cov}(Y(t), W(t)) &= \frac{t^2}{2}
		\end{align*}
	\end{frame}
	
	\begin{frame}{3.4: Exponential Integrated Brownian Motion}
		Another form of integrated Brownian motion arises if we assume the \textit{percent rate of change} of a price follows a Brownian motion process (similar to what we did in GBM).
		
		If $P(t)$ denotes the price at time $t$, then:
		$$\frac{d}{dt} P(t) = W(t)P(t)$$
		
		Solving this differential equation yields:
		$$P(t) = P(0) \exp\left\{ \int_0^t W(s) \, ds \right\} = P(0) e^{Y(t)}$$
		
		Taking $P(0) = 1$, we have $P(t) = e^{Y(t)}$. Since $Y(t)$ is normal with mean 0 and variance $t^3/3$, $P(t)$ is log-normally distributed.
		
		$$\mathbb{E}[P(t)] = \mathbb{E}\left[e^{Y(t)}\right] = \exp\left\{ \frac{t^3}{6} \right\}$$
		
		The above comes from the $M_Y(1)$ where $M_Y$ is the MGF of $Y(t)$.
	\end{frame}
	
	% ---------------------------------------------------
	% 3.5 Brownian Motion with Drift
	% ---------------------------------------------------
	\begin{frame}{3.5: Brownian Motion with Drift}
		A Brownian motion with drift represents a standard Brownian motion that is shifted by a constant trend over time.
		
		\vspace{0.5cm}
		\textbf{Definition:} The process $X(t)$ is defined as:
		$$X(t) = \mu t + \sigma W(t)$$
		
		\vspace{0.2cm}
		Where $\mu$ is the constant drift rate and $\sigma$ is the volatility. At time $t$, $X(t)$ is normally distributed with mean $\mu t$ and variance $\sigma^2 t$.
	\end{frame}
	
	\begin{frame}{3.5: Brownian Motion with Drift (Simulation)}
		\begin{figure}
			\centering
			\includegraphics[width=\textwidth]{"6_drift_brownian.png"}
			\caption{Paths and distribution of Brownian Motion with Drift}
		\end{figure}
	\end{frame}
	
	% ---------------------------------------------------
	% Slide 1: Problem Setup
	% ---------------------------------------------------
	\begin{frame}{3.5: Optimal Stopping Time}
		Let $X(t)$ be a Brownian motion with drift $\mu$ and variance $\sigma^2t$.
		
		\vspace{0.3cm}
		We want to compute the probability that the process hits an upper barrier $A > 0$ before a lower barrier $-B < 0$.
		
		\vspace{0.3cm}
		Let $P(x)$ denote this probability, conditioned on starting at a position $x \in (-B, A)$. So $X(0)$ need not be 0:
		$$P(x) = P(X(t) \text{ hits } A \text{ before } -B \mid X(0) = x)$$
		
		\vspace{0.3cm}
		To find $P(x)$, we will observe the process over a microscopic time interval $[0, h]$. Let $Y = X(h) - X(0)$ be the random step taken during this time.
	\end{frame}
	
	% ---------------------------------------------------
	% Slide 2: The Law of Total Probability
	% ---------------------------------------------------
	\begin{frame}{3.5: Optimal Stopping Time}
		Let $H$ be the event that the process hits either barrier ($A$ or $-B$) during the microscopic time interval $[0, h]$. Let Win be the event that the process hits $A$ before it hits $-B$.
		
		\vspace{0.3cm}
		By the Law of Total Probability, we can split $P(x)$ into two scenarios:
		\begin{align*}
			P(x) &= P(\text{Win} \mid H) P(H) \\
			&+ P(\text{Win} \mid H^c) P(H^c)
		\end{align*}
		
		\vspace{0.3cm}
		Because Brownian paths are continuous and $x$ is strictly inside $(-B, A)$, moving a macroscopic distance to a boundary in a infinitesimal time $h$ is small.
		
		\vspace{0.3cm}
		Mathematically, $P(H)$ shrinks to zero faster than $h$, so $P(H) = o(h)$.
	\end{frame}
	
	% ---------------------------------------------------
	% Slide 3: Isolating the Relevant Scenario
	% ---------------------------------------------------
	\begin{frame}{3.5: Optimal Stopping Time}
		Using $P(H) = o(h)$,
		\begin{align*}
			P(x) &= P(\text{Win} \mid H) \cdot o(h) + P(\text{Win} \mid H^c) (1 - o(h)) \\
			P(x) &= P(\text{Win} \mid H^c) + o(h)
		\end{align*}
		\textit{(Note that $o(h) \times bounded = o(h)$, here the difference of probabilities is bounded between $-1$ and $1$)}
		
		\vspace{0.3cm}
		This tells us that for an infinitesimally small step $h$, we can safely assume the process has not yet hit a boundary. 
		
	\end{frame}
	
	% ---------------------------------------------------
	% Slide 4: The Markov Property and Expectation
	% ---------------------------------------------------
	\begin{frame}{3.5: Optimal Stopping Time}
		Markov property dictates that the future evolution of the process is dependent only on the present:
		$$ \mathbb{P}(\text{Future} \mid \text{Present}, \text{Past}) = \mathbb{P}(\text{Future} \mid \text{Present}) $$
		$$ \mathbb{P}(W \mid X(h) = x+Y, H^C, X(0) = x) = \mathbb{P}(W \mid X(h) = x+Y) $$
		
		Because Brownian motion is \textbf{time-homogeneous}, the probability of winning starting from position $x+Y$ at time $h$ is identical to starting there at time $0$:
		$$ \mathbb{P}(W \mid X(h) = x+Y) = P(x+Y) $$
		
		Integrating over all possible values $Y=y$ of the random step $Y = X(h) - X(0)$ yields the total expectation: $\mathbb{E}[P(x+Y)]$.
	\end{frame}
	
	% ---------------------------------------------------
	% Slide 5: Taylor Series Expansion
	% ---------------------------------------------------
	\begin{frame}{3.5: Optimal Stopping Time}
		We now proceed by taking a Taylor series expansion of $P(x + Y)$ about $x$:
		$$P(x+Y) = P(x) + P'(x)Y + \frac{P''(x)}{2}Y^2 + \cdots$$
		
		\vspace{0.3cm}
		Substitute this expansion into our expectation equation:
		$$P(x) = \mathbb{E}\left[P(x) + P'(x)Y + \frac{P''(x)}{2}Y^2\right] + o(h)$$
		
		\vspace{0.3cm}
		Since $Y \sim \mathcal{N}(\mu h, h)$, we know its moments:
		\begin{itemize}
			\item $\mathbb{E}[Y] = \mu h$
			\item $\mathbb{E}[Y^2] = \text{Var}(Y) + (\mathbb{E}[Y])^2 = \sigma^2h + \mu^2 h^2 = h + o(h)$
		\end{itemize}
		\textit{(Higher-order terms like $Y^3$ yield expectations that are $o(h)$).}
	\end{frame}
	
	% ---------------------------------------------------
	% Slide 6: Forming the Differential Equation
	% ---------------------------------------------------
	\begin{frame}{3.5: Optimal Stopping Time}
		Applying the expectations to our expansion:
		$$P(x) = P(x) + P'(x)\mu h + \frac{P''(x)}{2}\sigma^2h + o(h)$$
		
		\vspace{0.3cm}
		Subtract $P(x)$ from both sides, divide the entire equation by $h$, and isolate the terms:
		$$P'(x)\mu + \frac{P''(x)}{2}\sigma^2 = \frac{o(h)}{h}$$
		
		\vspace{0.3cm}
		Letting the time step $h \to 0$, the $\frac{o(h)}{h}$ term vanishes by definition. This leaves us with a second-order linear differential equation:
		$$P'(x)\mu + \frac{P''(x)}{2}\sigma^2 = 0$$
	\end{frame}
	
	% ---------------------------------------------------
	% Slide 7: Solving the Differential Equation
	% ---------------------------------------------------
	\begin{frame}{3.5: Optimal Stopping Time}
		Integrating once yields $e^{2(\mu/\sigma^2) x} P'(x) = c_1$. 
		Integrating a second time gives the general solution:
		$$P(x) = C_1 + C_2 e^{-2(\mu/\sigma^2) x}$$
		
		We use the boundary conditions $P(A) = 1$ and $P(-B) = 0$ to solve for $C_1$ and $C_2$:
		\begin{align*}
			C_1 + C_2 e^{-2(\mu/\sigma^2) A} &= 1 \\
			C_1 + C_2 e^{2(\mu/\sigma^2) B} &= 0
		\end{align*}
		
		\vspace{0.2cm}
		Solving this linear system yields:
		$$C_1 = \frac{e^{2(\mu/\sigma^2) B}}{e^{2(\mu/\sigma^2) B} - e^{-2(\mu/\sigma^2) A}}, \quad C_2 = \frac{-1}{e^{2(\mu/\sigma^2) B} - e^{-2(\mu/\sigma^2) A}}$$
		
		\vspace{0.2cm}
		Substituting $C_1$ and $C_2$ back gives the probability from any starting point $x$:
		$$P(x) = \frac{e^{2(\mu/\sigma^2) B} - e^{-2(\mu/\sigma^2) x}}{e^{2(\mu/\sigma^2) B} - e^{-2(\mu/\sigma^2) A}}$$
	\end{frame}
	
\end{document}